% FIGURE NEEDED: circular diagram of the celestial equator seen from the
% north: Greenwich and Local Meridians, vernal equinox Y, and the arcs GST,
% GST*, LST, RA, L, LHA = 0 (2017_Theory_Solution.pdf, page 1).
When it culminates, $\mathrm{LHA} = 0$,
$\mathrm{RA} = -\mathrm{LST}$ \hfill[1.0]
\begin{align*}
\mathrm{LST} &= -(\mathrm{L} - \mathrm{GST}^{*}) = -\mathrm{RA}\\
\mathrm{GST}^{*} &= \mathrm{L} - \mathrm{RA}
  = 6^{\mathrm{h}}\,34^{\mathrm{min}} - 5^{\mathrm{h}}\,24^{\mathrm{min}}
  = 1^{\mathrm{h}}\,10^{\mathrm{min}}
\end{align*}
\hfill[2.0]
\begin{align*}
\mathrm{GST} &= 24^{\mathrm{h}} - 1^{\mathrm{h}}\,10^{\mathrm{min}}\\
\mathrm{GST} &= 22^{\mathrm{h}}\,50^{\mathrm{min}}
\end{align*}
\hfill[1.0]

Phuket time of 21:00 is UT\,14:00. \hfill[1.0]

\[ \mathrm{GST} = \text{GST of 1st Jan}
   + \text{The angle that $\Upsilon$ moves away from 1st Jan}
   + \mathrm{UT} \]
\hfill[1.0]
\begin{align*}
22^{\mathrm{h}}\,50^{\mathrm{min}}
  &= 6^{\mathrm{h}}\,43^{\mathrm{min}}
   + \left[\left(\frac{\text{Days that away from 1st Jan}}
                     {365.2422\ \text{days}}\right)\times 24\ \mathrm{h}\right]
   + \left[14^{\mathrm{h}}\times\left(\frac{24}{23.9344}\right)\right]\\
\text{Days that away from 1st Jan}
  &= (22^{\mathrm{h}}\,50^{\mathrm{min}} - 6^{\mathrm{h}}\,43^{\mathrm{min}}
    - 14^{\mathrm{h}}\,2^{\mathrm{min}})
    \left(\frac{365.2422\ \text{days}}{24^{\mathrm{h}}}\right)
\end{align*}
\hfill[2.0]

Days that away from 1st Jan $= 31.7$ days from 1st Jan (at Greenwich).
\hfill[1.0]

The time of the year that the LMC culminates in Phuket is on 2nd February.
\hfill[1.0]
